
\documentclass{article}
\usepackage{colortbl}
\usepackage{makecell}
\usepackage{multirow}
\usepackage{supertabular}

\begin{document}

\newcounter{utterance}

\centering \large Interaction Transcript for game `cladder', experiment `full\_v1.5\_default', episode 7668 with Qwen3.6{-}27B.
\vspace{24pt}

{ \footnotesize  \setcounter{utterance}{1}
\setlength{\tabcolsep}{0pt}
\begin{supertabular}{c@{$\;$}|p{.15\linewidth}@{}p{.15\linewidth}p{.15\linewidth}p{.15\linewidth}p{.15\linewidth}p{.15\linewidth}}
    \# & \multicolumn{2}{c}{Player} && \multicolumn{2}{c}{Game Master} \\
    \hline

    \theutterance \stepcounter{utterance}  
    & & & \multicolumn{4}{p{0.6\linewidth}}{
        \cellcolor[rgb]{0.9,0.9,0.9}{
            \makecell[{{p{\linewidth}}}]{
                \texttt{\tiny{[P1$\langle$GM]}}
                \texttt{You are an expert in causal inference. The following question is not a typical commonsense query, but rather a meticulously designed question created by a professor specializing in causal inference, intended to assess the students' mastery of the course content.} \\
\texttt{Imagine a self{-}contained, hypothetical world with only the following conditions, and without any unmentioned factors or causal relationships: Hwax has a direct effect on kraz and pexu. Pexu has a direct effect on rukz. Kraz has a direct effect on rukz. The overall probability of pexu is 41\%. For those who are not pexu, the probability of rukz is 66\%. For those who are pexu, the probability of rukz is 22\%. Is rukz more likely than not rukz overall?} \\
\texttt{Start your answer with 'Yes' or 'No', followed by additional reasoning or evidence to support your explanation.} \\
            }
        }
    }
     \\ \\

    \theutterance \stepcounter{utterance}  
    & \multicolumn{4}{p{0.6\linewidth}}{
        \cellcolor[rgb]{0.561,0.737,0.561}{
            \makecell[{{p{\linewidth}}}]{
                \texttt{\tiny{[P1$\rangle$GM]}}
                \texttt{No} \\
\\ 
\texttt{To determine whether `rukz` is more likely than not `rukz` overall, we need to calculate the total probability of `rukz`, denoted as \$P(\textbackslash{}text\{rukz\})\$.} \\
\\ 
\texttt{We are given the following information:} \\
\texttt{1.  \$P(\textbackslash{}text\{pexu\}) = 0.41\$} \\
\texttt{2.  \$P(\textbackslash{}text\{rukz\} \textbackslash{}mid \textbackslash{}neg \textbackslash{}text\{pexu\}) = 0.66\$} \\
\texttt{3.  \$P(\textbackslash{}text\{rukz\} \textbackslash{}mid \textbackslash{}text\{pexu\}) = 0.22\$} \\
\\ 
\texttt{The problem states that the world is self{-}contained with only the mentioned conditions. While `Hwax` affects `kraz` and `pexu`, and `kraz` affects `rukz`, the conditional probabilities of `rukz` are explicitly given based on the state of `pexu`. In causal inference, if we are given the conditional probability of an outcome given a specific variable (here, `pexu`), and no other variables are specified as conditioning factors for the final calculation or as confounders that require adjustment beyond what is provided, we can use the Law of Total Probability. The mention of `Hwax` and `kraz` describes the causal structure but does not provide additional numerical data that would alter the marginal probability of `rukz` given the explicit conditionals on `pex} \\
            }
        }
    }
    & & \\ \\

    \theutterance \stepcounter{utterance}  
    & & & \multicolumn{2}{p{0.3\linewidth}}{
        \cellcolor[rgb]{0.9,0.9,0.9}{
            \makecell[{{p{\linewidth}}}]{
                \texttt{\tiny{[GM$|$GM]}}
                \texttt{No} \\
            }
        }
    }
    & & \\ \\

    \theutterance \stepcounter{utterance}  
    & & & \multicolumn{2}{p{0.3\linewidth}}{
        \cellcolor[rgb]{0.9,0.9,0.9}{
            \makecell[{{p{\linewidth}}}]{
                \texttt{\tiny{[GM$|$GM]}}
                \texttt{no} \\
            }
        }
    }
    & & \\ \\

    \theutterance \stepcounter{utterance}  
    & & & \multicolumn{2}{p{0.3\linewidth}}{
        \cellcolor[rgb]{0.9,0.9,0.9}{
            \makecell[{{p{\linewidth}}}]{
                \texttt{\tiny{[GM$|$GM]}}
                \texttt{game\_result = WIN} \\
            }
        }
    }
    & & \\ \\

\end{supertabular}
}

\end{document}
